% Part: first-order-logic
% Chapter: beyond
% Section: intuitionistic-logic

\documentclass[../../../include/open-logic-section]{subfiles}

\begin{document}

\olfileid{fol}{byd}{il}

\olsection{Constructief redeneren met intuïtionistische logica}

Tweede-ordelogica en hogere-ordelogica voegen mogelijkheden toe aan de klassieke
eerste-ordelogica. Intuïtionistische eerste-ordelogica doet juist het omgekeerde:
zij staat een deel van het klassieke redeneren niet toe. Zo wil deze logica een
meer ``constructieve'' manier van redeneren weergeven. De volgende voorbeelden
laten zien wat daarmee wordt bedoeld.

Stel dat iemand op een dag zegt dat die een natuurlijk getal~$x$ heeft bepaald.
Het getal heeft deze eigenschap: als $x$ een priemgetal is, is de
Riemann-hypothese waar. Als $x$ samengesteld is, is de Riemann-hypothese onwaar.
Dat klinkt als geweldig nieuws. De vraag of de Riemann-hypothese waar is, is een
van de grote open vragen in de wiskunde. Deze persoon lijkt die vraag te hebben
teruggebracht tot een berekening: we hoeven alleen te bepalen of één vast getal
een priemgetal is.

Wat is dan de magische waarde van $x$? Het antwoord is: $x$ is het natuurlijke
getal dat gelijk is aan $7$ als de Riemann-hypothese waar is, en anders aan $9$.

Boos vraag je je geld terug. Volgens de klassieke logica legt deze beschrijving
inderdaad precies één waarde van $x$ vast. Maar jij wilt de waarde van $x$
\emph{expliciet} krijgen. Je wilt dus te horen krijgen welk getal het is.

Hier volgt een tweede voorbeeld dat misschien minder gekunsteld is. We weten dat
we een irrationaal getal tot een rationale macht kunnen verheffen en dan een
rationaal getal kunnen krijgen. Zo is $\sqrt{2}^2 = 2$. Maar kan een irrationaal
getal tot een \emph{irrationale} macht ook een rationaal resultaat geven? De
volgende stelling zegt van wel:

\begin{thm}
Er bestaan irrationale getallen $a$ en $b$ waarvoor $a^b$ rationaal is.
\end{thm}

\begin{proof}
Kijk naar $\sqrt{2}^{\sqrt{2}}$. Als dit getal rationaal is, zijn we klaar: neem
$a = b = \sqrt{2}$. Als het niet rationaal is, is het irrationaal. Dan geldt
\[
(\sqrt{2}^{\sqrt{2}})^{\sqrt{2}} = \sqrt{2}^{\sqrt{2} \cdot
  \sqrt{2}} = \sqrt{2}^2 = 2,
\]
en dit is zeker rationaal. Neem in dit tweede geval dus $a$ gelijk aan
$\sqrt{2}^{\sqrt{2}}$, en neem $b$ gelijk aan~$\sqrt 2$.
\end{proof}

Is dit een geldig bewijs? De meeste wiskundigen vinden van wel. Toch ontbreekt er
iets. We hebben bewezen dat er een paar reële getallen met de gevraagde
eigenschap bestaat, maar we kunnen niet zeggen \emph{welk} paar het is. Er is ook
een bewijs dat het paar $a$, $b$ \emph{wel} noemt: neem $a = \sqrt{3}$ en
$b = \log_3 4$. Dan
\[
a^b = \sqrt{3}^{\log_3 4} = 3^{1/2 \cdot \log_3 4} = (3^{\log_3
  4})^{1/2} = 4^{1/2}= 2,
\]
omdat $3^{\log_3 x} = x$.

Intuïtionistische logica is bedoeld voor een manier van redeneren waarin de stap
uit het eerste bewijs niet mag. Wil je bewijzen dat er een $x$ bestaat die
aan~$!A(x)$ voldoet, dan moet je een concrete~$x$ geven. Je moet er ook een bewijs
bij geven dat deze aan $!A$ voldoet, zoals in het tweede bewijs. Wil je bewijzen
dat $!A$ of $!B$ geldt, dan moet je een van die twee formules kunnen bewijzen.

Formeel krijgen we intuïtionistische eerste-ordelogica door !!a{derivation}
systeem voor eerste-ordelogica op een bepaalde manier te beperken. Er bestaan
ook intuïtionistische versies van tweede-ordelogica en hogere-ordelogica.
Wiskundig gezien zijn dit formele systemen met regels voor bewijzen. Ze zijn
tegelijk bedoeld als model van een bepaalde manier van wiskundig redeneren. Je
kunt vinden dat juist deze manier past bij een bepaalde kijk op de wiskunde,
zoals het intuïtionisme van Brouwer. Je kunt haar ook concreter en
bevredigender vinden, zoals in het constructivisme van Bishop. En je kunt erover
twisten of het formele systeem de informele bedoeling goed weergeeft. Welke
filosofische positie je ook kiest, we kunnen intuïtionistische logica gewoon als
een precies beschreven formele logica bestuderen. Veel wiskundige logici vinden
dat om allerlei redenen interessant.

Er is een informele, constructieve uitleg van de intuïtionistische connectieven.
Die heet de BHK-interpretatie, naar Brouwer, Heyting en Kolmogorov. De uitleg
zegt telkens wat je als bewijs moet leveren. Een bewijs van $!A \land !B$ bestaat
uit een bewijs van $!A$ samen met een bewijs van $!B$. Een bewijs van $!A \lor !B$
bestaat uit een bewijs van $!A$ of een bewijs van $!B$, met de uitdrukkelijke
informatie welke van de twee je hebt. Een bewijs van $!A \lif !B$ is een
procedure die elk bewijs van $!A$ omzet in een bewijs van~$!B$. Een bewijs van
$\lforall[x][!A(x)]$ is een procedure die een bewijs van $!A(x)$ teruggeeft voor
elke waarde van~$x$. Een bewijs van $\lexists[x][!A(x)]$ bestaat uit een waarde
van~$x$ en een bewijs dat die waarde aan~$!A$ voldoet. We kunnen dit ook in
rekentermen uitleggen. Dat heet het ``Curry--Howard-isomorfisme'' of het
``!!{formula}s-als-typenparadigma''. Zie !!a{formula} als een beschrijving van een
bepaald gegevenstype. Een bewijs is dan een rekenbaar object van dat type. Het
bestaan van zo'n object laat zien dat de bijbehorende !!{formula} waar is.

Intuïtionistische logica wordt vaak omschreven als klassieke logica ``zonder''
het principe van de uitgesloten derde. Dat principe zegt dat A of niet-A geldt.
De volgende stelling maakt deze omschrijving preciezer.
\begin{thm}
Binnen de intuïtionistische logica zijn de volgende axiomaschema's equivalent:
\begin{enumerate}
\item $(\lnot !A \lif \lfalse) \lif !A$.
\item $!A \lor \lnot !A$
\item $\lnot \lnot !A \lif !A$
\end{enumerate}
\end{thm}
Het is een goede oefening in intuïtionistische logica om uit elk van deze
schema's de gevallen van de andere twee af te leiden.

Kolmogorov, Glivenko en Heyting gaven onafhankelijk van elkaar de eerste formele
bewijssystemen voor de intuïtionistische propositielogica. Die systemen moesten
het intuïtionisme van Brouwer precies vastleggen. Heyting gaf de eerste
formalisering van de intuïtionistische eerste-ordelogica en van delen van de
intuïtionistische wiskunde. Sommige schema's die klassiek geldig zijn, zijn
intuïtionistisch niet geldig. Veel andere schema's zijn dat wel.

De \emph{dubbele-negatievertaling} laat een belangrijk verband zien tussen de
klassieke en de intuïtionistische logica. We definiëren de vertaling stap voor
stap volgens de opbouw van formules. Lees $!A^N$ als de ``intuïtionistische''
vertaling van de klassieke !!{formula}~$!A$:
\begin{align*}
!A^N & \ident \lnot\lnot !A \quad \text{voor atomaire !!{formula}s $!A$} \\
(!A \land !B)^N & \ident (!A^N \land !B^N) \\
(!A \lor !B)^N & \ident  \lnot\lnot (!A^N \lor !B^N) \\
(!A \lif !B)^N & \ident (!A^N \lif !B^N) \\
(\lforall[x][!A])^N & \ident \lforall[x][!A^N] \\
(\lexists[x][!A])^N & \ident \lnot\lnot\lexists[x][!A^N]
\end{align*}
Kolmogorov en Glivenko hadden versies van deze vertaling voor de
propositielogica. G\"odel en Gentzen gaven, onafhankelijk van elkaar, de
vertaling voor de predicatenlogica. Er geldt:

\begin{thm}
\begin{enumerate}
\item $!A \liff !A^N$ is klassiek bewijsbaar.
\item Als $!A$ klassiek bewijsbaar is, dan is $!A^N$ intuïtionistisch
  bewijsbaar.
\end{enumerate}
\end{thm}

Stel je nu het volgende gesprek voor. De klassieke wiskundige zegt: ``Ik heb
$!A$ bewezen!'' De intuïtionistische wiskundige antwoordt: ``Dat bewijs is niet
geldig. Je hebt eigenlijk $!A^N$ bewezen.'' De klassieke wiskundige zegt: ``Dat
vind ik ook goed!'' Voor de klassieke wiskundige is het verschil haarkloverij,
want de twee formules zijn klassiek equivalent. De intuïtionistische wiskundige
blijft erbij dat er wel een verschil is.

Deze vertaling gaat alleen over de zuivere logica. Zij zegt niet welke
\emph{niet-logische} axioma's we in de klassieke of de intuïtionistische
wiskunde moeten aannemen. Zij zegt ook niet hoe die twee groepen axioma's met
elkaar samenhangen. Een kleine uitbreiding van de vorige stelling geeft wel de
volgende nuttige informatie:

\begin{thm}
Als $\Gamma$ klassiek $!A$ bewijst, dan bewijst $\Gamma^N$
intuïtionistisch $!A^N$.
\end{thm}

Dus: als $!A$ klassiek uit bepaalde aannamen kan worden bewezen, dan kan
$!A^N$ uit de dubbele-negatievertalingen van die aannamen worden bewezen.

Om te laten zien dat een zin of propositionele !!{formula} intuïtionistisch
geldig is, hoef je alleen een bewijs te geven. Maar hoe laat je zien dat zij niet
geldig is? Daarvoor hebben we een semantiek nodig die correct is en liefst ook
volledig. De Kripke-semantiek werkt hiervoor goed.

We kunnen nu hetzelfde doen als bij de klassieke logica: eerst de semantiek
definiëren en daarna correctheid en volledigheid bewijzen. Er is wel een verschil.
Bij de klassieke logica leek de semantiek vanzelfsprekend. Zij leek al besloten
te liggen in de betekenis van de connectieven. We kunnen de Kripke-semantiek ook
intuïtief uitleggen, maar zij voelt minder onvermijdelijk. Bovendien is een
klassieke !!{structure} op zichzelf een natuurlijk wiskundig object. We kunnen
!!a{structure} gebruiken om klassieke eerste-ordelogica te bestuderen. We kunnen
ook andersom klassieke eerste-ordelogica gebruiken om wiskundige !!{structure}s
te bestuderen. Een Kripke-!!{structure} lijkt alleen een logisch hulpmiddel en
niet een wiskundig object dat los daarvan interessant is.

Een Kripke-!!{structure}~$\mModel M = \tuple{W, R, V}$ voor een propositionele
taal heeft drie onderdelen. Ten eerste is er een verzameling~$W$. Ten tweede is
er een partiële ordening~$R$ op~$W$ met een kleinste !!{element}. Ten derde
worden propositionele variabelen ``monotoon'' aan de !!{element}s van~$W$
toegekend. Monotoon betekent hier: wat in een toestand waar is, blijft in elke
latere toestand waar. De !!{element}s van $W$ stellen ``werelden'' of
``kennistoestanden'' voor. Een element $v \geq u$ stelt een ``mogelijke
toekomstige toestand'' van~$u$ voor. De propositionele variabelen die aan~$u$
zijn toegekend, zijn precies de proposities waarvan we in toestand~$u$ weten dat
ze waar zijn. De forceringsrelatie $\mSat{M}{!A}[w]$ breidt deze toekenning uit
naar alle !!{formula}s van de taal. Lees $\mSat{M}{!A}[w]$ als: ``$!A$ is waar
in toestand~$w$.'' We definiëren de relatie stap voor stap:
\begin{enumerate}
\item $\mSat{M}{\Obj p_i}[w]$ desda $\Obj p_i$ een van de propositionele
  variabelen is die aan~$w$ zijn toegekend.
\item $\mSat/{M}{\lfalse}[w]$.
\item $\mSat{M}{(!A \land !B)}[w]$ desda $\mSat{M}{!A}[w]$ en $\mSat{M}{!B}[w]$.
\item $\mSat{M}{(!A \lor !B)}[w]$ desda $\mSat{M}{!A}[w]$ of $\mSat{M}{!B}[w]$.
\item $\mSat{M}{(!A \lif !B)}[w]$ desda voor alle $w' \geq w$ geldt: als
  $\mSat{M}{!A}[w']$, dan $\mSat{M}{!B}[w']$.
\end{enumerate}
Probeer als oefening te laten zien dat $\lnot (p \land q) \lif
(\lnot p \lor \lnot q)$ niet intuïtionistisch geldig is. Maak daarvoor een
Kripke-!!{structure} waarin deze formule niet geldt. Die structuur is dan een
tegenvoorbeeld.

\end{document}
